% GENERATED FILE. Edit canonical lessons, then run npm run generate:books.
% canonical-source-hash: fnv1a64:baa477438e961ce8
% canonical-lessons: PA-C238-ata, PA-C238-parat, PA-C238-gunnhna, PA-C238-pera, PA-C238-velna, PA-R238-review-232-238

\chapter{Kneading the Dough}
\label{ch:pa-a2-kneading-the-dough}
\hlchaptermodality{238}

\begin{chapteropening}
\textbf{By the end of this chapter:} \emph{I can say how dough is made ready.}

Complete the last lesson of chapter 238: I can say how dough is made ready.
\end{chapteropening}

\section[\texorpdfstring{\pa{ਆਟਾ} (āṭā)}{ਆਟਾ (āṭā)}]{\pa{ਆਟਾ} (āṭā) — flour}

\label{lesson:PA-C238-ata}



\begin{quote}
Before the new one: say the Punjabi for a churning pot, then the Punjabi for a churning stick.
\end{quote}

\subsection*{You'll want to know: \pa{ਆਟਾ}}
\textbf{\pa{ਆਟਾ}} — \emph{āṭā} — \textquotedblleft{}flour\textquotedblright{}.

\begin{quote}
\pa{ਆਟਾ} \pa{ਡੱਬੇ} \pa{ਵਿੱਚ} \pa{ਹੈ।}
\end{quote}

\begin{quote}
\emph{āṭā ḍabbe vich hai.}
\end{quote}

\begin{quote}
The flour is in the box.
\end{quote}

A masculine noun.

\begin{cousinweb}
\textbf{\pa{ਆਟਾ}} (\emph{āṭā}) is inherited from a reconstructed \emph{*ārta}, \textquotedblleft{}flour\textquotedblright{}; CDIAL lists Punjabi \emph{āṭṭā}, \textquotedblleft{}flour\textquotedblright{} (CDIAL 1338).
\end{cousinweb}

\begin{grammarlens}[title={What you've built}]
One more word for this chapter.
\end{grammarlens}

\subsection*{Your turn}
\begin{itemize}
\raggedright
  \item \emph{Say it:} \emph{āṭā}
  \item \emph{Say it:} \emph{āṭā}, once more
  \item \emph{Say it:} \emph{āṭā ḍabbe vich hai.} — The flour is in the box.
  \item \emph{Recall:} say the Punjabi for a churning pot, then the Punjabi for a churning stick, then say \emph{āṭā} again
\end{itemize}

\begin{tcolorbox}[breakable,colback=teal!4,colframe=teal!35!black,title={Before you move on}]
What does \pa{ਆਟਾ} mean? (\textquotedblleft{}Flour\textquotedblright{}.) Say it once more.
\end{tcolorbox}
\section[\texorpdfstring{\pa{ਪਰਾਤ} (parāt)}{ਪਰਾਤ (parāt)}]{\pa{ਪਰਾਤ} (parāt) — a kneading tray}

\label{lesson:PA-C238-parat}



\begin{quote}
Before the new one: say the Punjabi for a churning stick, then the Punjabi for flour.
\end{quote}

\subsection*{You'll want to know: \pa{ਪਰਾਤ}}
\textbf{\pa{ਪਰਾਤ}} — \emph{parāt} — \textquotedblleft{}a kneading tray\textquotedblright{}.

\begin{quote}
\pa{ਪਰਾਤ} \pa{ਵਿੱਚ} \pa{ਆਟਾ} \pa{ਹੈ।}
\end{quote}

\begin{quote}
\emph{parāt vich āṭā hai.}
\end{quote}

\begin{quote}
There is flour in the kneading tray.
\end{quote}

A feminine noun.

\begin{cousinweb}
\textbf{\pa{ਪਰਾਤ}} (\emph{parāt}) is inherited from Sanskrit \emph{pātra}, \textquotedblleft{}drinking vessel, dish\textquotedblright{}; CDIAL lists Punjabi \emph{parāt}, \textquotedblleft{}large wooden kneading vessel\textquotedblright{} (CDIAL 8055).
\end{cousinweb}

\begin{grammarlens}[title={What you've built}]
One more word for this chapter.
\end{grammarlens}

\subsection*{Your turn}
\begin{itemize}
\raggedright
  \item \emph{Say it:} \emph{parāt}
  \item \emph{Say it:} \emph{parāt}, once more
  \item \emph{Say it:} \emph{parāt vich āṭā hai.} — There is flour in the kneading tray.
  \item \emph{Recall:} say the Punjabi for a churning stick, then the Punjabi for flour, then say \emph{parāt} again
  \item \emph{Write it:} copy \textbf{\pa{ਪਰਾਤ}} once, checking the whole word
\end{itemize}

\begin{tcolorbox}[breakable,colback=teal!4,colframe=teal!35!black,title={Before you move on}]
What does \pa{ਪਰਾਤ} mean? (\textquotedblleft{}A kneading tray\textquotedblright{}.) Say it once more.
\end{tcolorbox}
\section[\texorpdfstring{\pa{ਗੁੰਨ੍ਹਣਾ} (gunnhṇā)}{ਗੁੰਨ੍ਹਣਾ (gunnhṇā)}]{\pa{ਗੁੰਨ੍ਹਣਾ} (gunnhṇā) — to knead}

\label{lesson:PA-C238-gunnhna}



\begin{quote}
Before the new one: say the Punjabi for flour, then the Punjabi for a kneading tray.
\end{quote}

\subsection*{You'll want to know: \pa{ਗੁੰਨ੍ਹਣਾ}}
\textbf{\pa{ਗੁੰਨ੍ਹਣਾ}} — \emph{gunnhṇā} — \textquotedblleft{}to knead\textquotedblright{}.

\begin{quote}
\pa{ਮੈਂ} \pa{ਪਰਾਤ} \pa{ਵਿੱਚ} \pa{ਆਟਾ} \pa{ਗੁੰਨ੍ਹਦੀ} \pa{ਹਾਂ।}
\end{quote}

\begin{quote}
\emph{maiṁ parāt vich āṭā gunnhdī hāṁ.}
\end{quote}

\begin{quote}
I knead the flour in the tray.
\end{quote}

Past (I …): \pa{ਮੈਂ} \pa{ਆਟਾ} \pa{ਗੁੰਨ੍ਹਿਆ}, \emph{maiṁ āṭā gunnhiā}.

\begin{cousinweb}
\textbf{\pa{ਗੁੰਨ੍ਹਣਾ}} (\emph{gunnhṇā}) is inherited from Sanskrit \emph{guphati}, \textquotedblleft{}strings together\textquotedblright{}; CDIAL lists Punjabi \emph{gunnhṇā}, \textquotedblleft{}to braid, knead\textquotedblright{} (CDIAL 4205).
\end{cousinweb}

\begin{grammarlens}[title={What you've built}]
One more word for this chapter.
\end{grammarlens}

\subsection*{Your turn}
\begin{itemize}
\raggedright
  \item \emph{Say it:} \emph{gunnhṇā}
  \item \emph{Say it:} \emph{gunnhṇā}, once more
  \item \emph{Say it:} \emph{maiṁ parāt vich āṭā gunnhdī hāṁ.} — I knead the flour in the tray.
  \item \emph{Recall:} say the Punjabi for flour, then the Punjabi for a kneading tray, then say \emph{gunnhṇā} again
  \item \emph{Write it:} copy \textbf{\pa{ਗੁੰਨ੍ਹਣਾ}} once, checking the whole word
\end{itemize}

\begin{tcolorbox}[breakable,colback=teal!4,colframe=teal!35!black,title={Before you move on}]
What does \pa{ਗੁੰਨ੍ਹਣਾ} mean? (\textquotedblleft{}To knead\textquotedblright{}.) Say it once more.
\end{tcolorbox}
\section[\texorpdfstring{\pa{ਪੇੜਾ} (peṛā)}{ਪੇੜਾ (peṛā)}]{\pa{ਪੇੜਾ} (peṛā) — a ball of dough}

\label{lesson:PA-C238-pera}



\begin{quote}
Before the new one: say the Punjabi for a kneading tray, then the Punjabi for to knead.
\end{quote}

\subsection*{You'll want to know: \pa{ਪੇੜਾ}}
\textbf{\pa{ਪੇੜਾ}} — \emph{peṛā} — \textquotedblleft{}a ball of dough\textquotedblright{}.

\begin{quote}
\pa{ਮੈਂ} \pa{ਆਟੇ} \pa{ਦਾ} \pa{ਪੇੜਾ} \pa{ਬਣਾਉਂਦੀ} \pa{ਹਾਂ।}
\end{quote}

\begin{quote}
\emph{maiṁ āṭe dā peṛā baṇāundī hāṁ.}
\end{quote}

\begin{quote}
I make a ball of dough.
\end{quote}

A masculine noun.

\begin{cousinweb}
\textbf{\pa{ਪੇੜਾ}} (\emph{peṛā}) is inherited from a reconstructed \emph{*pēḍa}, \textquotedblleft{}lump\textquotedblright{}; CDIAL lists Punjabi \emph{peṛā}, \textquotedblleft{}lump\textquotedblright{} (CDIAL 8377).
\end{cousinweb}

\begin{grammarlens}[title={What you've built}]
One more word for this chapter.
\end{grammarlens}

\subsection*{Your turn}
\begin{itemize}
\raggedright
  \item \emph{Say it:} \emph{peṛā}
  \item \emph{Say it:} \emph{peṛā}, once more
  \item \emph{Say it:} \emph{maiṁ āṭe dā peṛā baṇāundī hāṁ.} — I make a ball of dough.
  \item \emph{Recall:} say the Punjabi for a kneading tray, then the Punjabi for to knead, then say \emph{peṛā} again
  \item \emph{Write it:} copy \textbf{\pa{ਪੇੜਾ}} once, checking the whole word
\end{itemize}

\begin{tcolorbox}[breakable,colback=teal!4,colframe=teal!35!black,title={Before you move on}]
What does \pa{ਪੇੜਾ} mean? (\textquotedblleft{}A ball of dough\textquotedblright{}.) Say it once more.
\end{tcolorbox}
\section[\texorpdfstring{\pa{ਵੇਲਣਾ} (velṇā)}{ਵੇਲਣਾ (velṇā)}]{\pa{ਵੇਲਣਾ} (velṇā) — to roll out (dough)}

\label{lesson:PA-C238-velna}



\begin{quote}
Before the new one: say the Punjabi for to knead, then the Punjabi for a ball of dough.
\end{quote}

\subsection*{You'll want to know: \pa{ਵੇਲਣਾ}}
\textbf{\pa{ਵੇਲਣਾ}} — \emph{velṇā} — \textquotedblleft{}to roll out (dough)\textquotedblright{}.

\begin{quote}
\pa{ਮੈਂ} \pa{ਪੇੜਾ} \pa{ਪਤਲਾ} \pa{ਵੇਲਦੀ} \pa{ਹਾਂ।}
\end{quote}

\begin{quote}
\emph{maiṁ peṛā patlā veldī hāṁ.}
\end{quote}

\begin{quote}
I roll the ball of dough out thin.
\end{quote}

Past (I …): \pa{ਮੈਂ} \pa{ਰੋਟੀ} \pa{ਵੇਲੀ}, \emph{maiṁ roṭī velī}.

\begin{cousinweb}
\textbf{\pa{ਵੇਲਣਾ}} (\emph{velṇā}) is inherited from Sanskrit \emph{vēllayati}, \textquotedblleft{}causes to shake, kneads\textquotedblright{}; CDIAL lists Punjabi \emph{velṇā}, \textquotedblleft{}to roll out\textquotedblright{} (CDIAL 12121).
\end{cousinweb}

\begin{grammarlens}[title={What you've built}]
One more word for this chapter.
\end{grammarlens}

\subsection*{Your turn}
\begin{itemize}
\raggedright
  \item \emph{Say it:} \emph{velṇā}
  \item \emph{Say it:} \emph{velṇā}, once more
  \item \emph{Say it:} \emph{maiṁ peṛā patlā veldī hāṁ.} — I roll the ball of dough out thin.
  \item \emph{Recall:} say the Punjabi for to knead, then the Punjabi for a ball of dough, then say \emph{velṇā} again
  \item \emph{Write it:} copy \textbf{\pa{ਵੇਲਣਾ}} once, checking the whole word
\end{itemize}

\begin{tcolorbox}[breakable,colback=teal!4,colframe=teal!35!black,title={Before you move on}]
What does \pa{ਵੇਲਣਾ} mean? (\textquotedblleft{}To roll out (dough)\textquotedblright{}.) Say it once more.
\end{tcolorbox}
\section[(dialogue)]{Looking back: from Putting Things Away to Kneading the Dough}

\label{lesson:PA-R238-review-232-238}



\begin{quote}
Say the Punjabi for to arrange neatly and to gather up.
\end{quote}

\subsection*{Your turn}
Say these aloud:

\begin{itemize}
\raggedright
  \item to arrange neatly, to gather up, memory, to come unstitched, to stuff
  \item rest, a cotton rug, to spread out, a thick cotton sheet, a fan
  \item sleep (noun), an oil lamp, to burn, to put out, a wooden bed
  \item to sweep, to dust, a floor cloth, mud plaster, to plaster
  \item to scrub, to rinse, to tip out, a small water pot, a deep bowl
  \item to boil, to boil down, yoghurt, a churning pot, a churning stick
  \item flour, a kneading tray, to knead, a ball of dough, to roll out
\end{itemize}

\begin{tcolorbox}[breakable,colback=teal!4,colframe=teal!35!black,title={Before you move on}]
To arrange neatly? (\textbf{\pa{ਸਜਾਉਣਾ}}, \emph{sajāuṇā}.) To gather up? (\textbf{\pa{ਸਮੇਟਣਾ}}, \emph{sameṭṇā}.) Memory? (\textbf{\pa{ਯਾਦ}}, \emph{yād}.) To come unstitched? (\textbf{\pa{ਉੱਧੜਨਾ}}, \emph{uddhaṛnā}.) To stuff? (\textbf{\pa{ਤੁੰਨਣਾ}}, \emph{tunnṇā}.) Rest? (\textbf{\pa{ਆਰਾਮ}}, \emph{ārām}.) A cotton rug? (\textbf{\pa{ਦਰੀ}}, \emph{darī}.) To spread out? (\textbf{\pa{ਵਿਛਾਉਣਾ}}, \emph{vichhāuṇā}.) A thick cotton sheet? (\textbf{\pa{ਖੇਸ}}, \emph{khes}.) A fan? (\textbf{\pa{ਪੱਖਾ}}, \emph{pakkhā}.) Sleep (noun)? (\textbf{\pa{ਨੀਂਦ}}, \emph{nīnd}.) An oil lamp? (\textbf{\pa{ਦੀਵਾ}}, \emph{dīvā}.) To burn? (\textbf{\pa{ਬਲਣਾ}}, \emph{balṇā}.) To put out? (\textbf{\pa{ਬੁਝਾਉਣਾ}}, \emph{bujhāuṇā}.) A wooden bed? (\textbf{\pa{ਪਲੰਘ}}, \emph{palaṅgh}.) To sweep? (\textbf{\pa{ਬੁਹਾਰਨਾ}}, \emph{buhārnā}.) To dust? (\textbf{\pa{ਝਾੜਨਾ}}, \emph{jhāṛnā}.) A floor cloth? (\textbf{\pa{ਪੋਚਾ}}, \emph{pochā}.) Mud plaster? (\textbf{\pa{ਗਾਰਾ}}, \emph{gārā}.) To plaster? (\textbf{\pa{ਲਿੱਪਣਾ}}, \emph{lippṇā}.) To scrub? (\textbf{\pa{ਮਾਂਜਣਾ}}, \emph{mānjṇā}.) To rinse? (\textbf{\pa{ਹੰਘਾਲਣਾ}}, \emph{haṅghālṇā}.) To tip out? (\textbf{\pa{ਉਲੱਦਣਾ}}, \emph{uladdṇā}.) A small water pot? (\textbf{\pa{ਲੋਟਾ}}, \emph{loṭā}.) A deep bowl? (\textbf{\pa{ਕਟੋਰਾ}}, \emph{kaṭorā}.) To boil? (\textbf{\pa{ਉਬਾਲਣਾ}}, \emph{ubālṇā}.) To boil down? (\textbf{\pa{ਔਟਣਾ}}, \emph{auṭṇā}.) Yoghurt? (\textbf{\pa{ਦਹੀਂ}}, \emph{dahīṁ}.) A churning pot? (\textbf{\pa{ਚਾਟੀ}}, \emph{chāṭī}.) A churning stick? (\textbf{\pa{ਮਧਾਣੀ}}, \emph{madhāṇī}.) Flour? (\textbf{\pa{ਆਟਾ}}, \emph{āṭā}.) A kneading tray? (\textbf{\pa{ਪਰਾਤ}}, \emph{parāt}.) To knead? (\textbf{\pa{ਗੁੰਨ੍ਹਣਾ}}, \emph{gunnhṇā}.) A ball of dough? (\textbf{\pa{ਪੇੜਾ}}, \emph{peṛā}.) To roll out? (\textbf{\pa{ਵੇਲਣਾ}}, \emph{velṇā}.)
\end{tcolorbox}
